\documentclass[11pt]{amsart}
\usepackage[T1]{fontenc}
\usepackage{lmodern}
\usepackage{amsmath,amssymb,amsthm,mathtools,mathrsfs}
\usepackage{microtype}
\usepackage[margin=1.05in]{geometry}
\usepackage{xcolor}
\usepackage[colorlinks=true,linkcolor=blue!50!black,citecolor=blue!50!black,urlcolor=blue!50!black]{hyperref}
\usepackage{enumitem,booktabs}
\setlist[enumerate]{leftmargin=2em,itemsep=2pt,topsep=3pt}
\setlist[itemize]{leftmargin=2em,itemsep=2pt,topsep=3pt}
\allowdisplaybreaks[2]
\newtheorem{theorem}{Theorem}[section]
\newtheorem{lemma}[theorem]{Lemma}
\newtheorem{proposition}[theorem]{Proposition}
\newtheorem{corollary}[theorem]{Corollary}
\theoremstyle{definition}
\newtheorem{definition}[theorem]{Definition}
\newtheorem{example}[theorem]{Example}
\theoremstyle{remark}
\newtheorem{remark}[theorem]{Remark}
\newcommand{\OO}{\mathcal O_F}
\newcommand{\kk}{\mathbb F_q}
\newcommand{\Gbar}{\mathrm{SL}_2(\kk)}
\newcommand{\G}{\mathrm{SL}_2(F)}
\newcommand{\wtG}{\widetilde G}
\newcommand{\wtK}{\widetilde K}
\newcommand{\cInd}{\operatorname{c\!-\!Ind}}
\newcommand{\tr}{\operatorname{tr}}
\newcommand{\vol}{\operatorname{vol}}
\newcommand{\Ad}{\operatorname{Ad}}
\newcommand{\Lie}{\operatorname{Lie}}
\newcommand{\vF}{v_F}
\title[Two-orbit depth-zero character relations]{Two-orbit unramified elliptic character relations for depth-zero representations of metaplectic covers of $\mathrm{SL}_2$}
\author{David Joyner}
\date{Working draft 4 --- November 11, 2026}
\subjclass[2020]{22E50, 22E35, 20G05}
\keywords{Metaplectic cover, depth-zero supercuspidal representation, Bruhat--Tits tree, elliptic torus, character identity, fixed subtree}
\begin{document}
\begin{abstract}
Let $F$ be a nonarchimedean local field of odd residue characteristic and let
$\widetilde G$ be an odd-degree tame central cover of $G=\mathrm{SL}_2(F)$.
We compare ordinary and genuine generic depth-zero supercuspidal
representations associated with both $G$-orbits of vertices in the
Bruhat--Tits tree.  The two-orbit theorem assumes specified homomorphic
splittings of the cover over representatives of both maximal compact
vertex stabilizers. Reparameterization of the finite nonsplit torus by
the inverse of the covering degree yields a pointwise character relation
on every regular element of either corresponding unramified elliptic
torus, including elements with central reduction.  A fixed-ball
calculation gives complementary explicit character formulas for the two
vertex orbits.  Their sum is simply the negative of the two torus
character values, with the depth-dependent terms cancelling.
The proof uses finite-group cuspidal character values, parity of the
tree, and uniqueness of splittings on pro-$p$ subgroups. Exceptional
finite types, ramified tori, positive-depth representations, and
orbital-integral transfer are not treated. No claim of novelty over the
existing character literature is made without a separate priority analysis.
\end{abstract}
\maketitle
\tableofcontents

\section{Introduction and the strengthened result}\label{sec:intro}
Let $F$ be a nonarchimedean local field, with ring of integers $\OO$, maximal ideal $\mathfrak p_F=(\varpi)$, residue field $\kk$ of odd cardinality $q$, and residue characteristic $p$.  Put
\[
 G=\G,\quad K=\mathrm{SL}_2(\OO),\quad K_1=\ker(K\to\Gbar).
\]
Consider a topological central extension
\begin{equation}\label{eq:cover}
 1\longrightarrow\mu_n\longrightarrow\wtG
 \xrightarrow{\mathbf p}G\longrightarrow1,
\end{equation}
where $n\ge3$ is odd, $(n,p)=1$, and $\mu_n\subset F^\times$.  We assume a \emph{fixed continuous homomorphic splitting}
\begin{equation}\label{eq:splitting}
 s_K:K\longrightarrow\wtG,\qquad \mathbf p\circ s_K=\mathrm{id}_K.
\end{equation}
The argument applies to any central extension with these properties.  For a particular metaplectic cocycle, the existence and normalization of $s_K$ must be verified separately; it is not implicit in the term ``metaplectic.''

Since $\mu_n\subset F^\times$ and $(n,p)=1$, the reduction of roots of unity gives
\begin{equation}\label{eq:coprime}
 n\mid(q-1),\qquad \gcd(n,q+1)=1.
\end{equation}
Let $E/F$ be the unramified quadratic extension.  Fix an $\OO$-basis of $\mathcal O_E$, and use it to identify $E$ with $F^2$ and $\mathcal O_E$ with $\OO^2$.  Multiplication by elements of $E$ then gives an integral embedding of its norm-one torus into $G$:
\begin{equation}\label{eq:torus}
 T=T(F)=\ker(N_{E/F}:E^\times\to F^\times)\subset K.
\end{equation}
The embedding is fixed throughout.  Its reduction is the nonsplit torus
\[
 \bar T=\ker(N_{\mathbb F_{q^2}/\kk}:\mathbb F_{q^2}^\times\to\kk^\times),
 \qquad |\bar T|=q+1.
\]
We identify the scalar elements $\pm1\in E^\times$ and their reductions with the matrices $\pm I$ in $G$ and $\Gbar$, respectively; both notations will be used for these same elements.  We write $T^{\mathrm{reg}}=T\setminus\{\pm I\}$.  These are regular elliptic elements in $G$; some reduce to $\pm I$.

Let $\theta:\bar T\to\mathbb C^\times$ be a character with $\theta^2\ne1$.  The corresponding finite-group cuspidal representation $\sigma_\theta$ of $\Gbar$ has dimension $q-1$.  If $a$ is an integer satisfying $an\equiv1\pmod{q+1}$, define
\begin{equation}\label{eq:theta-power}
 \theta^{\langle n\rangle}(u)=\theta(u^a)\qquad(u\in\bar T).
\end{equation}
Since $n$ is odd and $q+1$ is even, every inverse $a$ of $n$ modulo $q+1$ is odd.  Hence
\begin{equation}\label{eq:central-character}
 \theta^{\langle n\rangle}(-1)=\theta(-1).
\end{equation}
The regularity condition $\theta^2\ne1$ is preserved by this reparameterization.

Inflate $\sigma_\eta$ to $K$, and put
\[
 \pi_\eta=\operatorname{c\!-\!Ind}_K^G\sigma_\eta.
\]
Set $\wtK=\mathbf p^{-1}(K)=\mu_n s_K(K)$, choose the tautological character $\varepsilon:\mu_n\hookrightarrow\mathbb C^\times$, and define
\begin{equation}\label{eq:genuine-type}
 \widetilde\sigma_\eta(\zeta s_K(k))
 =\varepsilon(\zeta)\sigma_\eta(\bar k),\qquad
 \widetilde\pi_\eta=\operatorname{c\!-\!Ind}_{\wtK}^{\wtG}\widetilde\sigma_\eta.
\end{equation}
The standard compact-induction theorem and its lifted intertwining argument imply that these are irreducible depth-zero supercuspidal representations, with $\widetilde\pi_\eta$ genuine; see Proposition~\ref{prop:types} below.

\begin{theorem}[Unramified elliptic character identity]\label{thm:main}
With the preceding assumptions, let $\theta^2\ne1$ and $\eta=\theta^{\langle n\rangle}$.  Then for every $t\in T^{\mathrm{reg}}$,
\begin{equation}\label{eq:main}
 \boxed{\displaystyle
 \Theta_{\widetilde\pi_\eta}\bigl(s_K(t^n)\bigr)
 =\Theta_{\pi_\theta}(t).}
\end{equation}
Both sides are Harish--Chandra character functions on the regular set, with compatible normalizations of Haar measure.  The identity also holds after simultaneous conjugation by any element of $K$, using the homomorphism $s_K$ on the conjugated elements.
\end{theorem}

Theorem~\ref{thm:main} includes the regular-elliptic-reduction region of Draft~1 and, additionally, every element of $T$ congruent to $\pm I$ modulo $\mathfrak p_F$ except $\pm I$ itself. The latter elements have a nontrivial fixed subtree.  The new contribution of this draft is Theorem~\ref{thm:two-orbit}, treating compact induction from the \emph{other} vertex orbit as well.  The new hypothesis is a specified splitting over its stabilizer. The resulting complementary character formulas yield a striking cancellation in Corollary~\ref{cor:two-orbit-sum}. The character relation for all of $G_{\mathrm{reg}}$, ramified or positive-depth representations, and the fundamental lemma are not consequences of these theorems.

\begin{remark}[Why odd-degree covers, and not the double cover?]\label{rem:odd-only}
The hypothesis that $n$ is odd and prime to $q+1$ is essential to the present elementary construction, not merely convenient notation.  For $n=2$, the square map on the cyclic group $\bar T$ of even order $q+1$ is not bijective; thus the inverse exponent in \eqref{eq:theta-power} does not exist and the uniqueness assertion in Lemma~\ref{lem:power} fails.  Also, our proof that $s_K(-I)$ commutes with all of $\wtG$ uses the odd order of $\mu_n$.  Consequently the proof does not cover the twofold metaplectic group.  That case has its own Weil-representation and Shimura--Waldspurger structures and is related to the older correspondence of \cite{Joyner1999}; no negative claim about the existence of a separate twofold character correspondence is intended.
\end{remark}

\section{Finite-field character data}\label{sec:finite}
Write $\chi_\theta=\tr\sigma_\theta$.  We need only three rows of the character table of $\Gbar$.  By the signed Deligne--Lusztig formula and the Green function for the nonsplit torus \cite[\S2.3, Eq.~(2--6)]{CunninghamGordon}, the following standard statement holds.
\begin{proposition}[Three character values]\label{prop:finite}
For a regular $\theta$ of $\bar T$:
\begin{enumerate}[label=(\roman*)]
 \item If $u\in\bar T\setminus\{\pm1\}$, then
 \begin{equation}\label{eq:elliptic-finite}
  \chi_\theta(u)=-\theta(u)-\theta(u^{-1}).
 \end{equation}
 \item If $z\in\{\pm1\}$, then
 \begin{equation}\label{eq:central-finite}
  \chi_\theta(zI)=(q-1)\theta(z).
 \end{equation}
 \item If $z\in\{\pm1\}$ and $U\in\Gbar$ is nontrivial unipotent, then
 \begin{equation}\label{eq:unipotent-finite}
  \chi_\theta(zU)=-\theta(z).
 \end{equation}
 In particular, the value in (iii) is independent of the rational conjugacy class of $U$ and of $\theta$ except for its central character.
\end{enumerate}
\end{proposition}
\begin{proof}
For regular $\theta$, $\sigma_\theta=-R_{\bar T}^{\Gbar}(\theta)$.  The usual regular-semisimple formula gives (i).  The central element $zI$ acts by the scalar $\theta(z)$, proving (ii).  The Green function for the nonsplit torus takes value $+1$ on either nonidentity unipotent class; taking the negative of the Deligne--Lusztig character gives (iii) for $z=1$, and (ii)'s central scalar gives the case $z=-1$.  The exceptional quadratic parameter is excluded precisely because its constituents need not have the same values on the two rational unipotent classes.
\end{proof}

\begin{proposition}[Depth-zero compact induction]\label{prop:types}
For regular $\eta$, $\pi_\eta$ and $\widetilde\pi_\eta$ of \eqref{eq:genuine-type} are irreducible depth-zero supercuspidal representations of $G$ and $\widetilde G$, respectively.  The second is genuine, with central character $\varepsilon$ on $\mu_n$.
\end{proposition}
\begin{proof}
The classical compact-induction theorem applies to the inflation of the irreducible cuspidal $\sigma_\eta$ from the finite reductive quotient of $K$ \cite{MoyPrasad1996,CunninghamGordon}.  We give the additional intertwining argument for the cover.  Because $\widetilde K=\mu_n s_K(K)$, its quotient by $s_K(K_1)$ is $\mu_n\times\Gbar$ and the inducing type is a genuine twist of the ordinary inflation.

By the Cartan decomposition, every nonidentity double coset $K\backslash G/K$ has a representative
\[
 g_m=\begin{pmatrix}\varpi^m&0\\0&\varpi^{-m}\end{pmatrix},\qquad m\ge1.
\]
Put $H_m=K\cap g_mK_1g_m^{-1}$.  This is a pro-$p$ subgroup, although the \emph{larger} intersection $K\cap g_mKg_m^{-1}$ is generally not pro-$p$.  A direct matrix computation shows that reduction maps $H_m$ onto the lower unipotent radical
\[
 U^-(\kk)=\left\{\begin{pmatrix}1&0\\ c&1\end{pmatrix}:c\in\kk\right\}\subset\Gbar.
\]
For $h\in H_m$, let $u=g_m^{-1}hg_m\in K_1$.  Then $u\in K_1\cap g_m^{-1}Kg_m$.  Lemma~\ref{lem:no-phase}, applied with $x=g_m^{-1}$, gives
\[
 \widetilde g_m s_K(u)\widetilde g_m^{-1}=s_K(h),
 \quad\text{hence}\quad
 \widetilde g_m^{-1}s_K(h)\widetilde g_m=s_K(g_m^{-1}hg_m),
\]
for any chosen lift $\widetilde g_m$ of $g_m$.  The conjugated inducing representation is therefore trivial on $s_K(H_m)$, because $u\in K_1$.  If a nonzero intertwiner existed for this double coset, its image would lie in the $H_m$-fixed subspace of $\widetilde\sigma_\eta$.  But
\[
 (\widetilde\sigma_\eta)^{s_K(H_m)}
  =\sigma_\eta^{U^-(\kk)}=0,
\]
by cuspidality of $\sigma_\eta$.  Thus the intertwining space vanishes on every nonidentity double coset.  The Mackey irreducibility criterion yields the result for $\widetilde\pi_\eta$.  As $\widetilde K$ is compact modulo the finite center and the inducing type is finite-dimensional, the representation is supercuspidal.  Triviality on $s_K(K_1)$ gives depth zero, and genuineness follows from \eqref{eq:genuine-type}.
\end{proof}

\section{The torus and its fixed vertices}\label{sec:tree}
Let $\mathscr B$ be the Bruhat--Tits tree of $\mathrm{PGL}_2(F)$.  Its vertices are homothety classes of rank-two $\OO$-lattices in $F^2$.  Write $v_0=[\OO^2]$, whose stabilizer in $G$ is $K$.  The $G$-orbit $Gv_0$ comprises exactly the vertices at even distance from $v_0$.  The sphere of radius $d\ge1$ has
\begin{equation}\label{eq:spherecount}
 N_d=(q+1)q^{d-1}
\end{equation}
vertices.  These are elementary facts about the $(q+1)$-regular tree \cite{Casselman}.

\begin{lemma}[Power-map automorphism]\label{lem:power}
The map $t\mapsto t^n$ is an automorphism of the compact abelian group $T$.  Consequently, if $h\in G$ and $h^n=t^n$ for $t\in T^{\mathrm{reg}}$, then $h=t$.
\end{lemma}
\begin{proof}
The kernel $T_1$ of reduction $T\twoheadrightarrow\bar T$ is a pro-$p$ abelian group; the $n$-th power map is an automorphism on each finite $p$-group quotient because $(n,p)=1$, hence also on $T_1$.  The quotient $\bar T$ is cyclic of order $q+1$, on which the map is an automorphism by \eqref{eq:coprime}.  This proves the first claim.  The element $t^n$ is regular because $t^n=\pm I$ would imply $t=\pm I$, and its centralizer in $G$ is $T$.  Since $h$ commutes with $h^n=t^n$, it belongs to $T$; injectivity of the power map yields $h=t$.
\end{proof}

\begin{lemma}[Fixed vertex for regular reduction]\label{lem:singlevertex}
If $t\in T$ reduces to an element of $\bar T\setminus\{\pm1\}$, then its fixed-vertex set in $\mathscr B$ is $\{v_0\}$.  The same is true of $t^n$.
\end{lemma}
\begin{proof}
The reduction is regular elliptic and has no invariant line in $\kk^2$; thus $t$ fixes $v_0$ but no neighboring vertex.  An automorphism fixing two vertices fixes the entire geodesic between them, so it cannot fix any other vertex.  The power-map automorphism of $\bar T$ preserves $\bar T\setminus\{\pm1\}$ because $n$ is odd and prime to $q+1$.
\end{proof}

The new case begins with $t$ reducing to $zI$, where $z=\pm1$.  The element $z^{-1}t$ belongs to $K_1$.  Define its positive integral depth by
\begin{equation}\label{eq:r-def}
 r(t)=\max\{r\ge1: z^{-1}t\equiv I\pmod{\varpi^r}\}.
\end{equation}
This is finite for $t\ne zI$.

\begin{lemma}[Anisotropic first nonzero coefficient]\label{lem:anisotropic}
If $t\in T^{\mathrm{reg}}$ reduces to $zI$, write
\begin{equation}\label{eq:leading}
 t=z(I+\varpi^r A),\qquad r=r(t),\quad A\in M_2(\OO).
\end{equation}
Then $\bar A$ is a nonzero traceless matrix with irreducible characteristic polynomial over $\kk$.  In particular, for every $k\in K$ both off-diagonal entries of $\overline{k^{-1}Ak}$ are nonzero.
\end{lemma}
\begin{proof}
Choose a basis of the unramified quadratic algebra $\mathcal O_E$ in which $T$ acts by multiplication.  Its reduction is the field $\mathbb F_{q^2}$, and $\bar A\ne0$ because $r$ is maximal.  The determinant-one condition on $z^{-1}t$ implies $\tr\bar A=0$.  In $\mathbb F_{q^2}$, a nonzero trace-zero element has the form $b\sqrt\delta$ for some $b\in\kk^\times$ and a nonsquare $\delta\in\kk^\times$ (as $p\ne2$); its characteristic polynomial is $X^2-b^2\delta$, which is irreducible.  If either off-diagonal entry vanished in some basis over $\kk$, the corresponding coordinate line would be invariant, contradicting irreducibility.
\end{proof}

\begin{proposition}[Fixed-ball theorem]\label{prop:ball}
For $t$ as in Lemma~\ref{lem:anisotropic}, the fixed vertices of $t$ form the closed ball
\begin{equation}\label{eq:ball}
 \mathscr B^t=B(v_0,r(t))
  :=\{v\in\mathscr B:d(v,v_0)\le r(t)\}.
\end{equation}
Moreover, $r(t^n)=r(t)$ and $\mathscr B^{t^n}=\mathscr B^t$.
\end{proposition}
\begin{proof}
A vertex at distance $d$ from $v_0$ can be represented, after conjugation by a suitable $k\in K$, by the lattice
$D_d\OO^2$, where $D_d=\operatorname{diag}(\varpi^d,1)\in\mathrm{GL}_2(F)$.  Its stabilizer in $G$ is the intersection with $D_d\mathrm{GL}_2(\OO)D_d^{-1}$.  Writing $A'=k^{-1}Ak$, we have
\[
 D_d^{-1}(z^{-1}k^{-1}tk)D_d
 =I+\varpi^r
 \begin{pmatrix}
 a'_{11}&\varpi^{-d}a'_{12}\\
 \varpi^{d}a'_{21}&a'_{22}
 \end{pmatrix}.
\]
By Lemma~\ref{lem:anisotropic}, $a'_{12}$ is a unit.  The matrix lies in $\mathrm{SL}_2(\OO)$ exactly when $d\le r$, which proves \eqref{eq:ball}.  Since $n$ is prime to $p$ and odd,
\[
 t^n=z\bigl(I+n\varpi^r A+O(\varpi^{2r})\bigr),
\]
so its first nonzero coefficient modulo $\varpi$ is $n\bar A$, again nonzero and anisotropic.  Thus $r(t^n)=r(t)$ and the same ball description applies.
\end{proof}

\begin{remark}[Why parity matters]\label{rem:parity}
Although the fixed ball contains vertices at every distance $d\le r$, the coset space $G/K$ is indexed only by the \emph{even-distance} vertices.  Counting all vertices would produce an incorrect character formula.  In particular, a ball of radius one contributes just its center to the compact induction from $K$.  Concretely, for $d=2j$ the lattice homothety class represented by $D_{2j}=\operatorname{diag}(\varpi^{2j},1)\in\mathrm{GL}_2(F)$ is also represented by
\[
 \varpi^{-j}D_{2j}=\operatorname{diag}(\varpi^j,\varpi^{-j})\in G.
\]
The scalar factor has no effect on conjugation.  Consequently the calculation in Proposition~\ref{prop:ball} applies without change to the cosets used in the character sum below.
\end{remark}

\section{The central cover does not add phases on the relevant cosets}\label{sec:lift}
A splitting over $K$ does not generally commute with conjugation by elements outside $K$.  The next two lemmas explain why it does on the terms needed for Theorem~\ref{thm:main}.

\begin{lemma}[Unique conjugated splitting on a pro-$p$ subgroup]\label{lem:no-phase}
Let $x\in G$, choose any lift $\widetilde x\in\wtG$, and set
\[
 H=K_1\cap xKx^{-1}.
\]
For every $u\in H$,
\begin{equation}\label{eq:no-phase}
 \widetilde x^{-1}s_K(u)\widetilde x
  =s_K(x^{-1}ux).
\end{equation}
\end{lemma}
\begin{proof}
Both sides of \eqref{eq:no-phase}, as functions of $u\in H$, are continuous homomorphisms $H\to\wtG$ lifting the same homomorphism $u\mapsto x^{-1}ux$ into $G$.  Their ratio is therefore a continuous homomorphism $H\to\mu_n$.  Since $H$ is a closed subgroup of the pro-$p$ group $K_1$ and $(n,p)=1$, every such homomorphism is trivial.
\end{proof}

\begin{lemma}[Lifting the center]\label{lem:center-lift}
The element $s_K(-I)$ is central in $\wtG$.  Consequently, if $t=z u$ with $z\in\{\pm I\}$ and $u\in K_1$, then for every $x\in G$ such that $x^{-1}tx\in K$,
\begin{equation}\label{eq:lift-compatible}
 \widetilde x^{-1}s_K(t)\widetilde x
  =s_K(x^{-1}tx).
\end{equation}
\end{lemma}
\begin{proof}
The image of $s_K(-I)$ is central in $G$, so its commutator with any $\widetilde x$ lies in $\mu_n$.  Because $s_K(-I)$ has order two, this central commutator also has order dividing two.  As $n$ is odd, the commutator is trivial.  This proves the first statement.  For the second, $x^{-1}tx\in K$ is equivalent to $x^{-1}ux\in K$; apply Lemma~\ref{lem:no-phase} to $u$, then multiply by the central lift of $z$.
\end{proof}

\begin{remark}[Anti-genuine test functions]
The convention $\widetilde\pi(\zeta)=\varepsilon(\zeta)\mathrm{Id}$ requires \emph{anti-genuine} test functions $f(\zeta\widetilde g)=\varepsilon(\zeta)^{-1}f(\widetilde g)$ in the distribution-character pairing.  The pointwise identities here refer to the character function evaluated at specified lifts and use that convention.
\end{remark}

\section{Compact-induction sums and proof of the main theorem}\label{sec:proof}
Normalize Haar measure so that $\vol(K)=1$ and independently $\vol(\wtK)=1$.  The standard compact-induction formula at a regular elliptic element $t\in K$ is a finite sum
\begin{equation}\label{eq:ordinary-sum}
 \Theta_{\pi_\theta}(t)
 =\sum_{\substack{xK\in G/K\\ x^{-1}tx\in K}}
  \chi_\theta(\overline{x^{-1}tx}).
\end{equation}
Here right cosets correspond to vertices in $Gv_0$, and the summation is finite because the elliptic regular element has a finite fixed subtree.  For the genuine compact induction, the analogous sum uses the character of $\widetilde\sigma_\eta$ at $\widetilde x^{-1}s_K(t)\widetilde x$.

\begin{lemma}[Comparison at the same group element]\label{lem:cover-ordinary}
For every regular character $\eta$ of $\bar T$ and every $t\in T^{\mathrm{reg}}$,
\begin{equation}\label{eq:same-element}
 \Theta_{\widetilde\pi_\eta}(s_K(t))
 =\Theta_{\pi_\eta}(t).
\end{equation}
\end{lemma}
\begin{proof}
If $\bar t\notin\{\pm I\}$, Lemma~\ref{lem:singlevertex} says the unique fixed coset is $K$, and both sides are $\chi_\eta(\bar t)$.  If $\bar t=zI$, then $t=zu$ with $u\in K_1$.  Each fixed coset $xK$ contributes the ordinary trace $\chi_\eta(\overline{x^{-1}tx})$, because the lift of its conjugate is exactly $s_K(x^{-1}tx)$ by Lemma~\ref{lem:center-lift}.  The two finite sums therefore agree term by term.
\end{proof}

\begin{lemma}[Ordinary inverse-power comparison]\label{lem:ordinary-comparison}
For every $t\in T^{\mathrm{reg}}$ and $\eta=\theta^{\langle n\rangle}$,
\begin{equation}\label{eq:ordinary-comparison}
 \Theta_{\pi_\eta}(t^n)=\Theta_{\pi_\theta}(t).
\end{equation}
\end{lemma}
\begin{proof}
If $\bar t\notin\{\pm I\}$, Lemma~\ref{lem:singlevertex} reduces both sums to one term.  Equations \eqref{eq:theta-power} and \eqref{eq:elliptic-finite} give
\[
 \chi_\eta(\bar t^{\,n})
 =-\eta(\bar t^{\,n})-\eta(\bar t^{-n})
 =-\theta(\bar t)-\theta(\bar t^{-1})
 =\chi_\theta(\bar t).
\]
Now suppose $\bar t=zI$.  By Proposition~\ref{prop:ball}, $t$ and $t^n$ fix precisely the same vertices, hence the two sums in \eqref{eq:ordinary-sum} have the same index set.  For every contributing $xK$, set
\[
 y=\overline{x^{-1}tx}\in\Gbar.
\]
Its characteristic polynomial is $(X-z)^2$, so it is either $zI$ or $zU$ for a nontrivial unipotent $U$.  As $p\nmid n$ and $n$ is odd, $y^n$ is respectively $zI$ or $zU^n$, with $U^n$ again nontrivial unipotent.  By Proposition~\ref{prop:finite} and \eqref{eq:central-character},
\[
 \chi_\eta(y^n)=\chi_\theta(y)
\]
in either case.  Thus the character sums agree term by term.
\end{proof}

\begin{proof}[Proof of Theorem~\ref{thm:main}]
Apply Lemma~\ref{lem:cover-ordinary} to $t^n$ and then Lemma~\ref{lem:ordinary-comparison}:
\[
 \Theta_{\widetilde\pi_\eta}(s_K(t^n))
 =\Theta_{\pi_\eta}(t^n)
 =\Theta_{\pi_\theta}(t).
\]
The compact-induced representations have the asserted properties by Proposition~\ref{prop:types}.  The statement for $K$-conjugates follows because $s_K$ is a group homomorphism and Harish--Chandra characters are conjugation invariant.
\end{proof}

\section{Closed character formulas near the center}\label{sec:explicit}
The proof gives more than an equality: it provides an elementary character formula on the entire chosen unramified torus.

\begin{theorem}[Explicit fixed-ball character values]\label{thm:values}
Suppose $t\in T^{\mathrm{reg}}$ has reduction $zI$, $z\in\{\pm1\}$, and $r=r(t)\ge1$.  For every regular $\theta$,
\begin{equation}\label{eq:explicit}
 \boxed{\displaystyle
 \Theta_{\pi_\theta}(t)
 =\theta(z)\begin{cases}
 q^r-1,&r\text{ odd},\\
 -q^r-1,&r\text{ even}.
 \end{cases}}
\end{equation}
Equivalently, this is the single formula
\[
 \Theta_{\pi_\theta}(t)
 =\theta(z)\bigl((-1)^{r+1}q^r-1\bigr).
\]
The same value is obtained for $\Theta_{\widetilde\pi_\theta}(s_K(t))$.  The formula is independent of $\theta$ apart from the central sign $\theta(z)$.
\end{theorem}
\begin{proof}
By Proposition~\ref{prop:ball}, the fixed vertices are the ball of radius $r$, but only even-distance vertices belong to $G/K$.  For a fixed vertex at distance $d<r$, the integral matrix displayed in the proof of Proposition~\ref{prop:ball} reduces to $zI$, so the contributing trace is $(q-1)\theta(z)$ by \eqref{eq:central-finite}.  If $r$ is even, there are also fixed vertices at distance $d=r$; the reduction at each such vertex is $z$ times a nontrivial unipotent matrix, giving contribution $-\theta(z)$ by \eqref{eq:unipotent-finite}.

If $r=2m+1$, every even $d\le r$ satisfies $d<r$, and by \eqref{eq:spherecount}
\[
 \Theta_{\pi_\theta}(t)
 =\theta(z)(q-1)\left[1+\sum_{j=1}^{m}(q+1)q^{2j-1}\right]
 =\theta(z)(q^{2m+1}-1).
\]
If $r=2m$, the boundary sphere contributes negatively:
\[
 \begin{aligned}
 \Theta_{\pi_\theta}(t)
 &=\theta(z)\left\{(q-1)\left[1+\sum_{j=1}^{m-1}(q+1)q^{2j-1}\right]
 -(q+1)q^{2m-1}\right\}\\
 &=\theta(z)(-q^{2m}-1).
 \end{aligned}
\]
The genuine statement follows from Lemma~\ref{lem:cover-ordinary}.
\end{proof}

\begin{remark}[Comparison with Sally--Shalika]\label{rem:sally}
The formula at $z=I$ agrees with the Sally--Shalika expression as reproduced in \cite[Thm.~5.2]{AubertPlymen}: near the identity, the ordinary generic depth-zero character is $-1-f(t)$, and for elements in our unramified norm-one torus with $r(t)=r$ the factor is $f(t)=(-1)^r q^r$.  Multiplication by $z=-I$ changes the character by the central sign $\theta(-1)$.  This comparison supplies an external check on both the alternating sign and the constant term $-1$.
\end{remark}

\begin{remark}[Different roles of the parameter]\label{rem:two-regimes}
When $\bar t\in\bar T\setminus\{\pm I\}$, the reparameterization
$\eta=\theta^{\langle n\rangle}$ is used explicitly in the character formula
$-\eta(\bar t^n)-\eta(\bar t^{-n})$.  By contrast, when $\bar t=zI$, the value depends on $\theta$ only through $\theta(z)$; the comparison in this region needs only $r(t^n)=r(t)$ and $\eta(-1)=\theta(-1)$.
\end{remark}

\begin{corollary}[Restricted root-sum form]\label{cor:root-sum}
For $t\in T^{\mathrm{reg}}$ let
\[
 \Delta(t)=\left|\det(1-\Ad(t)\mid\mathfrak{sl}_2(F)/\Lie(T))\right|_F^{1/2}.
\]
Then $\Delta(t^n)=\Delta(t)$ and
\begin{equation}\label{eq:root-sum}
 \Delta(t^n)\Theta_{\widetilde\pi_{\theta^{\langle n\rangle}}}
  \bigl(s_K(t^n)\bigr)
 =\sum_{\substack{h\in G\\h^n=t^n}}
 \Delta(h)\Theta_{\pi_\theta}(h).
\end{equation}
The sum has exactly one term, $h=t$.
\end{corollary}
\begin{proof}
The uniqueness of $h$ is Lemma~\ref{lem:power}.  For reduction $\bar t\notin\{\pm1\}$, both discriminants are $1$.  In the central-reduction case $t=z(I+\varpi^r A)$ with anisotropic $\bar A$, the two nontrivial root directions have eigenvalues whose distances from $1$ have valuation $r$.  Thus $\Delta(t)=q^{-r}$, and the same argument for $t^n$ gives $\Delta(t^n)=q^{-r}$, since $(n,p)=1$.  Equation \eqref{eq:root-sum} follows from Theorem~\ref{thm:main}.
\end{proof}

\begin{remark}[Scope of the root-sum relation]
The lift in \eqref{eq:root-sum} is specifically $s_K(t^n)$.  Multiplication by $\zeta\in\mu_n$ multiplies the genuine character by $\varepsilon(\zeta)$.  No transfer factor for arbitrary lifts, no equality for ramified elliptic classes, and no matching orbital-integral identity is established.
\end{remark}


\section{The second vertex orbit and two-orbit character identities}\label{sec:two-orbit}
The compact subgroup $K$ in the preceding sections stabilizes the
vertex $v_0$.  The action of $G$ on the tree has \emph{two} vertex orbits;
the other orbit is needed to describe the full generic depth-zero family.
This section adds it without assuming an automorphism of $\widetilde G$
extending conjugation by $\mathrm{GL}_2(F)$.

Put
\begin{equation}\label{eq:second-parahoric}
 D=\begin{pmatrix}\varpi&0\\0&1\end{pmatrix}\in\mathrm{GL}_2(F),
 \quad v_1=Dv_0,\quad K^{(0)}=K,\quad K^{(1)}=DKD^{-1}.
\end{equation}
For $i\in\{0,1\}$ put
\[
 v_i=D^iv_0,\qquad K^{(i)}_+=D^iK_1D^{-i},\qquad
 T^{(i)}=D^iTD^{-i}.
\]
Here $K^{(i)}_+$ denotes the pro-$p$ radical of the \emph{vertex}
stabilizer, not another maximal compact subgroup.  The groups $K^{(0)}$
and $K^{(1)}$ are not $G$-conjugate.  Their quotients by the indicated
pro-$p$ radicals are identified with $\Gbar$ via
\[
 \rho_i:K^{(i)}\longrightarrow\Gbar,\qquad
 \rho_i(k)=\overline{D^{-i}kD^i}.
\]
The orbit $Gv_i$ consists of vertices whose distance from $v_0$ has
parity $i$.  Conjugation by $D$ acts on the tree and interchanges the
parities, although $D\notin G$.

\begin{definition}[Splittings and types for both orbits]\label{def:two-splits}
In addition to the splitting $s_0=s_K$ from \eqref{eq:splitting},
\emph{assume} a continuous homomorphic splitting
\begin{equation}\label{eq:second-splitting}
 s_1:K^{(1)}\longrightarrow\widetilde G,
 \qquad\mathbf p\circ s_1=\mathrm{id}_{K^{(1)}}.
\end{equation}
For a regular character $\eta$ of $\bar T$, let $\sigma_\eta^{(i)}$
be the inflation of $\sigma_\eta$ along $\rho_i$ and put
\begin{align}\label{eq:two-types}
 \pi_{i,\eta}&=\cInd_{K^{(i)}}^G\sigma_\eta^{(i)},\\
 \widetilde\pi_{i,\eta}
 &=\cInd_{\mathbf p^{-1}(K^{(i)})}^{\widetilde G}
 \widetilde\sigma_{i,\eta},\qquad
 \widetilde\sigma_{i,\eta}(\zeta s_i(k))
 =\varepsilon(\zeta)\sigma_\eta(\rho_i(k)).\nonumber
\end{align}
Thus $\pi_{0,\eta}=\pi_\eta$ and
$\widetilde\pi_{0,\eta}=\widetilde\pi_\eta$ in the notation of
Sections~\ref{sec:intro}--\ref{sec:explicit}.
\end{definition}

\begin{remark}[A genuine additional hypothesis]\label{rem:second-split}
The assumption \eqref{eq:second-splitting} must be checked for the
chosen cocycle.  Conjugation by $D$ is an automorphism of $G$, but it
is \emph{not} conjugation by an element of $G$; a lift of this outer
automorphism to an arbitrary central extension is not automatic.
We therefore do not manufacture $s_1$ by conjugating $s_0$ in
$\widetilde G$.  Once both splittings are given, they need not be
related globally on the intersection of the two maximal compact
subgroups: the proof requires compatibility only on pro-$p$
intersections and on $\{\pm I\}$, where it follows from uniqueness.
This formulation applies, in particular, to any tame metaplectic
realization for which the two compact splittings have been supplied.
\end{remark}

The following is the precise generalization of
Lemmas~\ref{lem:no-phase} and \ref{lem:center-lift} used for the
second orbit.
\begin{lemma}[Mixed pro-$p$ compatibility]\label{lem:mixed}
For $i,j\in\{0,1\}$, $x\in G$, and a lift
$\widetilde x\in\widetilde G$ of $x$, set
\[
 H_{i,j}(x)=K^{(i)}_+\cap xK^{(j)}x^{-1}.
\]
For $u\in H_{i,j}(x)$ one has
\begin{equation}\label{eq:mixed}
 \widetilde x^{-1}s_i(u)\widetilde x
 =s_j(x^{-1}ux).
\end{equation}
Moreover, $s_0(-I)=s_1(-I)$, and this common lift is central in
$\widetilde G$.  Consequently \eqref{eq:mixed} also holds with $u$
replaced by $zu$ for $z\in\{\pm I\}$ whenever $zu\in K^{(i)}$
and $x^{-1}(zu)x\in K^{(j)}$.
\end{lemma}
\begin{proof}
The group $H_{i,j}(x)$ is a closed subgroup of the pro-$p$ group
$K^{(i)}_+$.  On it, the two sides of \eqref{eq:mixed} are continuous
homomorphisms lifting the same map into $G$; their ratio is a
continuous homomorphism to $\mu_n$ and therefore trivial because
$(n,p)=1$.  Both $s_i(-I)$ have order two and project to $-I$.
Their ratio belongs to the odd-order central group $\mu_n$;
therefore the lifts coincide.  The commutator of this lift with
any element of $\widetilde G$ belongs to $\mu_n$ and has order
at most two, so it vanishes. Multiplication by this central lift
proves the last assertion.
\end{proof}

\begin{lemma}[Reduction at a fixed vertex]\label{lem:vertex-reduction}
Let $i\in\{0,1\}$ and $t\in (T^{(i)})^{\mathrm{reg}}$.
If $\rho_i(t)\notin\{\pm I\}$, the fixed set of $t$ is $\{v_i\}$.
If $\rho_i(t)=zI$, define $r_i(t)\ge1$ by
\[
 D^{-i}tD^i=z(I+\varpi^{r_i(t)} A),\qquad
 \bar A\ne0.
\]
Then $\mathscr B^t=B(v_i,r_i(t))$.  At any fixed vertex $v$ at
distance $d$ from $v_i$, the reduction of $t$ in the quotient of
its stabilizer is $zI$ if $d<r_i(t)$ and $zU$ for a nonidentity
unipotent $U$ if $d=r_i(t)$.  Also $r_i(t^n)=r_i(t)$ and
$\mathscr B^{t^n}=\mathscr B^t$.
\end{lemma}
\begin{proof}
Conjugation by $D^{-i}$ identifies this problem with
Lemmas~\ref{lem:singlevertex} and \ref{lem:anisotropic} and
Proposition~\ref{prop:ball}.  For the reduction statement,
use the displayed matrix in the proof of Proposition~\ref{prop:ball}.
When $d<r$, every off-diagonal entry has positive valuation and
the diagonal entries reduce to $z$.  When $d=r$, precisely the
upper-right entry may be a unit and is nonzero modulo $\varpi$;
the reduction is $z$ times a nontrivial unipotent element.
The resulting conjugacy type is independent of the chosen
vertex representative.  The final statement follows from
$p\nmid n$ and the anisotropic leading coefficient.
\end{proof}

\begin{theorem}[Two vertex orbits]\label{thm:two-orbit}
Assume both compact splittings \eqref{eq:splitting} and
\eqref{eq:second-splitting}.  Fix a regular character
$\theta$ of $\bar T$ and put $\eta=\theta^{\langle n\rangle}$.
For each pair $i,j\in\{0,1\}$ and every
$t\in (T^{(i)})^{\mathrm{reg}}$ one has
\begin{equation}\label{eq:two-orbit-identity}
 \boxed{\displaystyle
 \Theta_{\widetilde\pi_{j,\eta}}(s_i(t^n))
 =\Theta_{\pi_{j,\theta}}(t).}
\end{equation}
Thus the identity holds for both inducing vertex orbits when
evaluated on either of the two associated unramified elliptic
tori, including the cross-orbit terms.
\end{theorem}
\begin{proof}
The standard depth-zero compact-induction argument proving
Proposition~\ref{prop:types} applies to $K^{(1)}$ as well:
conjugation by $D$ identifies the ordinary intertwining
calculation with that for $K$, while uniqueness of pro-$p$
splittings handles the genuine intertwining comparison.  Hence
the four families in \eqref{eq:two-types} are irreducible
depth-zero supercuspidal representations.

For either $j$, the compact-induction character at
$t\in T^{(i)}$ is a finite sum indexed by the fixed vertices
belonging to $Gv_j$.  If $t$ has regular elliptic reduction
at $v_i$, Lemma~\ref{lem:vertex-reduction} says that only $v_i$
is fixed.  When $i\ne j$ there are no contributing vertices;
both sides of \eqref{eq:two-orbit-identity} are zero.  When
$i=j$, the sum has one term and
\[
 \chi_\eta(\rho_i(t)^n)=\chi_\theta(\rho_i(t))
\]
by \eqref{eq:elliptic-finite} and the inverse-power
reparameterization \eqref{eq:theta-power}.  In this single-coset
case the lift is given by $s_i$ itself.

If $t=z u$ with $u\in K^{(i)}_+$, the fixed balls of $t$
and $t^n$ agree by Lemma~\ref{lem:vertex-reduction}; thus the
ordinary character sums for $\pi_{j,\theta}(t)$ and
$\pi_{j,\eta}(t^n)$ have the same index set.  At each fixed
vertex the reduction is $zI$ or $zU$ with $U\ne I$
unipotent.  Its $n$-th power remains of the same kind;
by Proposition~\ref{prop:finite} and
$\eta(-1)=\theta(-1)$, the corresponding finite character
values are equal.  Hence
\[
 \Theta_{\pi_{j,\eta}}(t^n)=\Theta_{\pi_{j,\theta}}(t).
\]
For a contributing coset $xK^{(j)}$, Lemma~\ref{lem:mixed}
identifies the conjugate of the chosen lift $s_i(t^n)$
with $s_j(x^{-1}t^nx)$.  The genuine compact-induction
sum therefore equals the ordinary sum for
$\pi_{j,\eta}(t^n)$. Combining these equalities proves
\eqref{eq:two-orbit-identity}.
\end{proof}

\begin{theorem}[Complementary fixed-ball formulas]\label{thm:complementary}
Let $i,j\in\{0,1\}$ and
$t\in(T^{(i)})^{\mathrm{reg}}$.  Write
$\bar t_i=\rho_i(t)$.  If $\bar t_i\notin\{\pm I\}$, then
\begin{equation}\label{eq:two-orbit-far}
 \Theta_{\pi_{j,\theta}}(t)=
 \begin{cases}
 -\theta(\bar t_i)-\theta(\bar t_i^{-1}),&j=i,\\
 0,&j\ne i.
 \end{cases}
\end{equation}
If $\bar t_i=zI$ and $r=r_i(t)\ge1$, then
\begin{equation}\label{eq:two-orbit-near}
 \boxed{\displaystyle
 \Theta_{\pi_{j,\theta}}(t)
 =\theta(z)\left((-1)^{r+1+j-i}q^r-1\right).}
\end{equation}
Here only the parity of $j-i$ is relevant.  The corresponding
genuine character at $s_i(t)$ has the same value for the
unmodified parameter $\theta$.
\end{theorem}
\begin{proof}
The first assertion follows immediately from the unique-fixed-
vertex part of Lemma~\ref{lem:vertex-reduction}.  For the second,
only vertices at distance $d\le r$ from $v_i$ contribute; $Gv_j$
is precisely the subset of vertices for which
$d\equiv j-i\pmod 2$.  Each vertex with $d<r$ contributes
$(q-1)\theta(z)$, while a vertex with $d=r$ contributes
$-\theta(z)$.  The case $j=i$ is Theorem~\ref{thm:values},
transferred by conjugation with $D^i$.

For $j\ne i$, only \emph{odd} distances occur.  If $r=2m$,
all these vertices are interior, and the sphere count
\eqref{eq:spherecount} gives
\[
 \begin{aligned}
 \Theta_{\pi_{j,\theta}}(t)
  &=(q-1)\theta(z)\sum_{k=0}^{m-1}(q+1)q^{2k}\\
  &=\theta(z)(q^{2m}-1).
 \end{aligned}
\]
If $r=2m+1$, the last odd sphere is the boundary, so
\[
 \begin{aligned}
 \Theta_{\pi_{j,\theta}}(t)
  &=\theta(z)\left((q-1)\sum_{k=0}^{m-1}
     (q+1)q^{2k}-(q+1)q^{2m}\right)\\
  &=\theta(z)(-q^{2m+1}-1).
 \end{aligned}
\]
These are the claimed complementary signs.
The genuine identity at $s_i(t)$ follows from the mixed-splitting
argument in Theorem~\ref{thm:two-orbit} without taking powers.
\end{proof}

\begin{corollary}[The two-orbit sum]\label{cor:two-orbit-sum}
For every $i\in\{0,1\}$ and
$t\in (T^{(i)})^{\mathrm{reg}}$, put
$u=\rho_i(t)\in\bar T$.  Then
\begin{equation}\label{eq:packet-sum}
 \boxed{\displaystyle
 \Theta_{\pi_{0,\theta}}(t)+\Theta_{\pi_{1,\theta}}(t)
 =-\theta(u)-\theta(u^{-1}).}
\end{equation}
In particular, when $u=zI$ this sum is simply $-2\theta(z)$,
\emph{independent of the depth} $r_i(t)$.
For $\eta=\theta^{\langle n\rangle}$ the matching sum of genuine
characters satisfies
\begin{equation}\label{eq:genuine-packet-sum}
 \sum_{j=0}^1\Theta_{\widetilde\pi_{j,\eta}}(s_i(t^n))
 =-\theta(u)-\theta(u^{-1}).
\end{equation}
\end{corollary}
\begin{proof}
For regular elliptic reduction, only the $j=i$ summand is
nonzero, and it is the first line of
\eqref{eq:two-orbit-far}.  For central reduction, the
$q^r$ terms of \eqref{eq:two-orbit-near} cancel, leaving
$-2\theta(z)=-\theta(z)-\theta(z^{-1})$.  The last assertion
is the sum of the two identities \eqref{eq:two-orbit-identity}.
\end{proof}

\begin{remark}[Scope of the sum]
Formula \eqref{eq:packet-sum} is a concrete sum of two specific
depth-zero characters.  Although its cancellation resembles the
stable-sum formulas in the depth-zero $L$-packet literature
\cite{AubertPlymen}, the statement here does not identify the pair
with an $L$-packet, establish stability on all regular semisimple
conjugacy classes, or construct an endoscopic transfer factor.
The covering-group assertion is likewise confined to the specified
lifts of the two unramified tori.
\end{remark}

\section{Example over \texorpdfstring{$\mathbb Q_7$}{Q7}}\label{sec:example}
Take $F=\mathbb Q_7$, $n=3$, and $E=\mathbb Q_7(\sqrt3)$; $3$ is a nonsquare modulo $7$, so $E/F$ is unramified.  Realize $E$ on the basis $(1,\sqrt3)$ by multiplication:
\[
 a+b\sqrt3\longmapsto
 \begin{pmatrix}a&3b\\b&a\end{pmatrix}.
\]
The reduced norm-one torus is cyclic of order $8$.  Choose a generator $u\in\bar T$ and $\theta(u)=e^{2\pi i/8}$.  Since $3^{-1}\equiv3\pmod8$, the parameter associated to the genuine representation is $\eta=\theta^{\langle3\rangle}=\theta^3$.

For a lift $t_0\in T$ of $u$, both $t_0$ and $t_0^3$ have regular elliptic reduction; the character comparison yields
\[
 \Theta_{\pi_\theta}(t_0)
 =-2\cos(\pi/4)=-\sqrt2
 =\Theta_{\widetilde\pi_\eta}(s_K(t_0^3)).
\]
For a new, nontrivial test of the extension, define
\begin{equation}\label{eq:tr}
 t_r=\frac{1+7^r\sqrt3}{1-7^r\sqrt3}\in T,\qquad r\ge1.
\end{equation}
These elements are regular, reduce to $I$, and satisfy $r(t_r)=r$.  Their $7$-adic matrix coefficients have leading term $2\cdot7^r\sqrt3$, so Proposition~\ref{prop:ball} applies.  The following values follow from Theorem~\ref{thm:values}:
\begin{center}
\begin{tabular}{@{}cccc@{}}
\toprule
$r$ & even-distance fixed vertices in $G/K$ & $\Theta_{\pi_\theta}(t_r)$ & $\Theta_{\widetilde\pi_\eta}(s_K(t_r^3))$\\
\midrule
$1$ & $1$ & $6$ & $6$\\
$2$ & $1+56=57$ & $-50$ & $-50$\\
$3$ & $1+56=57$ & $342$ & $342$\\
\bottomrule
\end{tabular}
\end{center}
For $r=2$, the center vertex contributes $\chi_\theta(I)=6$, whereas the $56$ boundary vertices at distance two contribute $-1$ each, yielding $6-56=-50$.  For $r=3$, those same $56$ vertices lie strictly inside the fixed ball, so each contributes $6$; hence $57\cdot6=342$.  Raising $t_r$ to its third power preserves $r$ and the fixed ball.  This example illustrates exactly what is gained beyond the one-vertex theorem.

For the \emph{other} vertex orbit, the same elements produce new
values. The two orbits contribute complementary terms:
\begin{center}
\begin{tabular}{@{}rrrr@{}}
\toprule
$r$ & $\Theta_{\pi_{0,\theta}}(t_r)$
& $\Theta_{\pi_{1,\theta}}(t_r)$ & sum\\
\midrule
$1$ & $6$ & $-8$ & $-2$\\
$2$ & $-50$ & $48$ & $-2$\\
$3$ & $342$ & $-344$ & $-2$\\
\bottomrule
\end{tabular}
\end{center}
For $r=1$ there are eight boundary vertices in the odd orbit,
each contributing $-1$. For $r=2$ they lie in the interior and
each contributes $6$. For $r=3$ the eight first-sphere vertices
contribute $48$, while the $392$ boundary vertices contribute
$-392$. This gives $48-392=-344$ and verifies the new formula
in Theorem~\ref{thm:complementary} in three cases. All six
individual values satisfy the simultaneous metaplectic relation
in Theorem~\ref{thm:two-orbit}, conditional on the second
compact splitting.

\section{Relation to the literature and next problems}\label{sec:limits}
The ordinary depth-zero character calculations in the present paper are elementary specializations of the well-developed theory of characters of $\mathrm{SL}_2(F)$.  Sally--Shalika computed its supercuspidal character tables; Cunningham--Gordon gave motivic formulas for depth-zero characters on topologically unipotent elements; and Aubert--Plymen revisited the Sally--Shalika formulas in the setting of endoscopic character identities \cite{SallyShalika,CunninghamGordon,AubertPlymen}.  The finite Green-function values used here are standard.  The purpose of this reconstruction is to put the covering-group comparison, including the suppression of extra cocycle factors on pro-$p$ intersections and the two vertex orbits, into a direct proof accessible from the earlier manuscript \cite{Joyner1999}. The equalities of two-orbit ordinary characters are consistent with the depth-zero $L$-packet character formulas discussed by Aubert--Plymen; no stronger novelty claim is inferred.

Several issues remain outside this result:
\begin{enumerate}[label=(\roman*)]
\item \emph{Verifying compact splittings for specified cocycles.}  Theorem~\ref{thm:two-orbit} assumes a splitting over the second maximal compact.  Existence and normalization should be documented for each intended metaplectic realization.
\item \emph{Exceptional depth-zero types.}  The constituents associated with the quadratic character of $\bar T$ have different values on the two rational unipotent classes.  The simple termwise power comparison must therefore be reconsidered.
\item \emph{Ramified elliptic tori and positive depth.}  Their fixed-point geometry and character formulas are different; the present fixed-ball theorem does not apply.
\item \emph{A full character transfer or fundamental lemma.}  A pointwise identity on a particular elliptic torus and for a restricted family of representations does not establish either a representation-theoretic equivalence for the full genuine category or an orbital-integral transfer.
\end{enumerate}

\paragraph{Possible next theorem.}
A useful next target is a carefully normalized comparison of the exceptional depth-zero packet \emph{as a sum of characters}.  Such a formulation may avoid the obstruction from the separate rational unipotent classes. Another natural task is proving, for a concrete Kubota-type cocycle, the existence of both prescribed compact splittings and relating their normalizations.  Only after these tests should one attempt the ramified and positive-depth cases.

\section*{Acknowledgments and provenance}
This is a corrected and extended reconstruction of ideas in David Joyner's earlier unfinished manuscript \emph{Character relations for supercuspidal representations of $\overline{SL(2)}$}, represented by the surviving files \texttt{CHART.TEX}, \texttt{CHAR0.TEX}, \texttt{CHAR1.TEX}, and \texttt{CHAR2.TEX}.  Draft~1 restricted the identity to elements with regular elliptic reduction; Draft~2 included central-reduction elements and proved an explicit fixed-ball character formula.  Draft~3 clarified the lifted Mackey intertwining argument, the integral torus embedding, the parity convention, and the odd-degree limitation. Draft~4 extends the construction to the second vertex orbit, proves cross-orbit character comparisons under an explicit second-splitting hypothesis, and derives the complementary character formula and depth-independent two-orbit sum. The technical observations in an independent referee report provided by Argon-Gemini informed the earlier corrections.  The mathematical review, revision, and exposition were prepared with assistance from ChatGPT (OpenAI).  No new-priority claim is made without a full comparison with the existing character literature.

\begin{thebibliography}{99}
\bibitem{AubertPlymen}
A.-M.~Aubert and R.~Plymen,
\emph{Endoscopic character identities for depth-zero supercuspidal representations of $\mathrm{SL}_2$},
arXiv:2410.20183 (2024).
\url{https://arxiv.org/abs/2410.20183}.

\bibitem{Casselman}
W.~Casselman, \emph{The Bruhat--Tits tree of $\mathrm{SL}_2$}, notes (2014).
\url{https://personal.math.ubc.ca/~cass/research/pdf/Tree.pdf}.

\bibitem{CunninghamGordon}
C.~Cunningham and J.~Gordon,
\emph{Motivic proof of a character formula for $\mathrm{SL}(2)$},
Experiment. Math. \textbf{18} (2009), no.~1, 11--44.
\url{https://doi.org/10.1080/10586458.2009.10128882}.

\bibitem{DeligneLusztig}
P.~Deligne and G.~Lusztig,
\emph{Representations of reductive groups over finite fields},
Ann. of Math. (2) \textbf{103} (1976), 103--161.

\bibitem{Joyner1999}
D.~Joyner,
\emph{A correspondence for the supercuspidal representations of the metaplectic cover of $SL(2,F)$, $F$ $p$-adic},
Arch. Math. (Basel) \textbf{73} (1999), 332--340.
\url{https://doi.org/10.1007/s000130050406}.

\bibitem{MoyPrasad1996}
A.~Moy and G.~Prasad,
\emph{Jacquet functors and unrefined minimal $K$-types},
Comment. Math. Helv. \textbf{71} (1996), 98--121.

\bibitem{SallyShalika}
P.~J.~Sally, Jr. and J.~A.~Shalika,
\emph{Characters of the discrete series of representations of $SL(2)$ over a local field},
Proc. Nat. Acad. Sci. U.S.A. \textbf{61} (1968), 1231--1237.
\end{thebibliography}
\end{document}
